% Station 1 "Adaptation" - Deep Learning for Computer Vision, HSBI, WS 2026/27
% Template: CVPR 2026 author kit, https://github.com/cvpr-org/author-kit (tag CVPR2026-v1(latex)).
% Structure and page budget follow the module requirements (lecture 01, slide 16):
%   Abstract 100-150 words | Introduction 1/2-3/4 p. | Related Work 3/4-1 p. (>= 5 scientific sources)
%   Method 1-1 1/2 p. | Results 1-1 1/2 p. | Discussion 1-1 1/2 p.
%   max. 6 pages; references and appendix do not count.
% Outline with all placeholders: docs/milestones/MS1.md. Numbers only from results/summary.md.
% Build: make paper   (or: cd paper/ms1 && latexmk -pdf main.tex)

\documentclass[10pt,twocolumn,letterpaper]{article}

%%%%%%%%% PAPER TYPE  (author kit: final = camera-ready, author names, no ruler, no page numbers)
\usepackage{cvpr}              % To produce the CAMERA-READY version (submission to the module)
% \usepackage[review]{cvpr}      % To produce the REVIEW version (anonymous, ruler)
% \usepackage[pagenumbers]{cvpr} % To force page numbers, e.g. for drafts

%% This file contains a number of tweaks that are typically applied to the main document.
%% They are not enabled by default, but can be enabled by uncommenting the relevant lines.

%%
%% Inline annotations; for predefined colors, refer to "dvipsnames" in the xcolor package:
%% https://tinyurl.com/overleaf-colors
%%
\newcommand{\red}[1]{{\color{red}#1}}
\newcommand{\todo}[1]{{\color{red}#1}}
\newcommand{\TODO}[1]{\textbf{\color{red}[TODO: #1]}}
%%
%% disable for camera ready / submission by uncommenting these lines  
%%
% \renewcommand{\TODO}[1]{}
% \renewcommand{\todo}[1]{#1}

%%
%% work harder in optimizing text layout. Typically shrinks text by 1/6 of page, enable
%% it at the very end of the writing process, when you are just above the page limit
%%
% \usepackage{microtype}

%%
%% fine-tune paragraph spacing
%%
% \renewcommand{\paragraph}[1]{\vspace{.5em}\noindent\textbf{#1.}}

%%
%% globally adjusts space between figure and caption
%%
% \setlength{\abovecaptionskip}{.5em}


%%
%% Allows "the use of \paper to refer to the project name"
%% with automatic management of space at the end of the word
%%
% \usepackage{xspace}
% \newcommand{\paper}{ProjectName\xspace}

%%
%% Commonly used math definitions
%%
% \DeclareMathOperator*{\argmin}{arg\,min}
% \DeclareMathOperator*{\argmax}{arg\,max}

%%
%% Tigthen underline
%%
% \usepackage{soul}
% \setuldepth{foobar}
\usepackage{needspace}
% Columns that consist only of floats: figures start at the top instead of floating with stretchable space above.
\makeatletter
\setlength{\@fptop}{0pt}
\setlength{\@fpsep}{12pt plus 1fil}
\setlength{\@fpbot}{0pt plus 1fil}
\makeatother


% Results produced by `make station1` are included automatically once they exist.
\newcommand{\resultsdir}{../../results}
\newcommand{\resultfigure}[3]{% #1 file below results/, #2 width, #3 TODO text
  \IfFileExists{\resultsdir/#1}{\includegraphics[width=#2]{\resultsdir/#1}}{\fbox{\parbox{0.9\linewidth}{\TODO{#3 (\texttt{results/\detokenize{#1}})}}}}}
% Same, but limited by height (tall figures).
\newcommand{\resultfigureh}[3]{%
  \IfFileExists{\resultsdir/#1}{\includegraphics[height=#2]{\resultsdir/#1}}{\fbox{\parbox{0.9\linewidth}{\TODO{#3 (\texttt{results/\detokenize{#1}})}}}}}
% Matriculation numbers: local file matrikel.tex (git-ignored, defines \matrline); optional.
\IfFileExists{matrikel.tex}{\input{matrikel}}{\newcommand{\matrline}{}}
% Rows of the method table (decided here, because \IfFileExists cannot start a table cell).
\makeatletter
\IfFileExists{\resultsdir/ms1/tables.tex}%
  {\def\methodrows{\@@input \resultsdir/ms1/tables.tex }}%
  {\def\methodrows{\multicolumn{5}{c}{\TODO{run make station1 -- rows from results/ms1/tables.tex}}\\}}
\makeatother

% German paper (allowed by the module): babel for hyphenation and "Abbildung/Tabelle/Literatur",
% German names for \cref (cleveref is loaded by cvpr.sty at the end of the preamble; this hook runs later).
\usepackage[ngerman]{babel}
\AtEndPreamble{%
  \crefname{section}{Abschn.}{Abschn.}\Crefname{section}{Abschnitt}{Abschnitte}%
  \crefname{figure}{Abb.}{Abb.}\Crefname{figure}{Abbildung}{Abbildungen}%
  \crefname{table}{Tab.}{Tab.}\Crefname{table}{Tabelle}{Tabellen}%
  \crefname{equation}{Gl.}{Gl.}\Crefname{equation}{Gleichung}{Gleichungen}%
  \crefname{appendix}{Anhang}{Anhang}\Crefname{appendix}{Anhang}{Anhang}%
  \AtBeginDocument{\def\crefpairconjunction{ und }\def\crefmiddleconjunction{, }%
    \def\creflastconjunction{ und }\def\crefrangeconjunction{ bis }}%
}

\definecolor{cvprblue}{rgb}{0.21,0.49,0.74}
\usepackage[breaklinks,colorlinks,allcolors=cvprblue]{hyperref}

\def\paperID{MS1}
\def\confName{CVPR}
\def\confYear{2026}

%%%%%%%%% TITLE
\title{Wie effizient lässt sich ein vortrainierter Vision Transformer anpassen?\\Linear Probing, LoRA und CNNs auf 25 Hunderassen}

%%%%%%%%% AUTHORS (same institution: one block, see author kit)
\author{Wael Eskeif \qquad Tarek Dawalibi\\
{\small\matrline}\\
Hochschule Bielefeld (HSBI), Campus Minden\\
{\small Deep Learning for Computer Vision, WS 2026/27}\\
{\small Code: \url{https://github.com/hsbi-dlcv-eskeif-dawalibi/hunderassen-vit} (Tag \texttt{ms1})}
}

\begin{document}
\raggedbottom
\maketitle
% Abstract: 100-150 Wörter (Folie 16). Problem - Vorgehen - Kernergebnis.
\begin{abstract}
Vortrainierte Vision Transformer (ViTs) lassen sich gut auf neue Aufgaben übertragen, doch alle
Parameter anzupassen ist teuer. Wir untersuchen, wie effizient sich ein vortrainierter ViT-B/16 an die
Klassifikation der 25 Hunderassen von Oxford-IIIT Pet anpassen lässt. Mit einheitlichem Protokoll
und drei Seeds vergleichen wir Linear Probing, Low-Rank Adaptation (LoRA) mit verschiedenen Rängen und
Ziel-Layern, volles Fine-Tuning und zwei CNNs. LoRA mit Rang 8 erreicht die
höchste Test-Accuracy (96,2\,\%; die auf Validation gewählte Variante $r=4$ 96,1\,\%) bei 0,36\,\% trainierbaren Parametern, doch Linear Probing mit nur
0,02\,\% ist innerhalb der Seed-Streuung gleich gut. Volles Fine-Tuning liegt 0,8 Prozentpunkte darunter,
das beste CNN 1,4 Punkte unter dem besten ViT. Zusatzexperimente liefern Hinweise auf die hohe Genauigkeit: 22 der 25
Rassen sind ImageNet-Klassen, und die ursprünglichen ImageNet-Köpfe erkennen sie ganz ohne Training zu
96--97\,\% (Auswahl unter diesen 22). Bei Daten so nah am Pretraining ist ein neuer linearer Kopf die
effizienteste Anpassung.
\end{abstract}

% Einleitung: 1/2-3/4 Seite, mit konkreten Beiträgen (Folie 16).
\section{Einleitung}
\label{sec:intro}

Auf großen Datensätzen vortrainierte Vision Transformer (ViTs)~\cite{dosovitskiy2021vit} sind starke
Ausgangspunkte für neue Bildklassifikationsaufgaben. Bei kleinen Zieldatensätzen stellt sich die Frage,
\emph{wie} man sie anpasst: Nur einen neuen Klassifikator auf eingefrorenen Merkmalen zu trainieren
(Linear Probing) ist günstig, trennt die Zielklassen aber womöglich nicht gut genug. Alle Gewichte zu
aktualisieren (volles Fine-Tuning) ist teuer und kann die vortrainierten Merkmale
verzerren~\cite{kumar2022finetuning}. Low-Rank Adaptation (LoRA)~\cite{hu2022lora} bietet einen
Mittelweg: Die vortrainierten Gewichte bleiben eingefroren, gelernt werden nur kleine Änderungen niedrigen
Rangs.

Wir untersuchen diesen Zielkonflikt an einer feinkörnigen Aufgabe. Der Empfehlung folgend, nur eine
Tierart zu verwenden, klassifizieren wir die 25 Hunderassen von Oxford-IIIT Pet~\cite{parkhi2012pets};
Rassen derselben Art sehen sich ähnlich und unterscheiden sich in Details von Kopfform, Fell und
Körperproportionen. Die Situation ist typisch für die Praxis: wenige gelabelte Bilder (etwa 90 pro
Rasse für das Training), viele ähnliche Klassen und begrenzte Rechenleistung.

\paragraph{Forschungsfrage.}
Wie effizient lässt sich ein vortrainierter ViT an eine neue Aufgabe anpassen -- welche Accuracy
erreichen Linear Probing und LoRA im Vergleich zu vollem Fine-Tuning und zu CNNs, und zu welchen Kosten?
Wir prüfen vier Hypothesen:
\textbf{A1} LoRA mit $<1\,\%$ trainierbaren Parametern liegt weniger als 2 Prozentpunkte unter vollem
Fine-Tuning;
\textbf{A2} Linear Probing ist bei feinkörnigen Rassen deutlich schlechter als LoRA;
\textbf{A3} ein höherer LoRA-Rang als $r=8$ bringt wenig, mehr angepasste Layer bringen mehr;
\textbf{A4} der beste ViT und die CNNs erreichen eine ähnliche Accuracy.

\paragraph{Beiträge.}
\begin{itemize}
  \item Ein Vergleich von Linear Probing, LoRA und vollem Fine-Tuning des ViT-B/16 mit ResNet-50 und
        ConvNeXt-T unter identischem Protokoll (drei Seeds, Modellauswahl nur auf Validation).
  \item Eine LoRA-Ablation über Rang und angepasste Layer.
  \item Eine Kosten-Nutzen-Analyse: Accuracy gegenüber trainierbaren Parametern und Trainingszeit sowie
        der Einfluss der Datenaugmentierung.
  \item Eine Analyse, \emph{warum} alle Methoden über 90\,\% erreichen (Training ohne Pretraining,
        Überschneidung mit den ImageNet-Klassen, Suche nach Fast-Duplikaten), und Erklärungskarten,
        worauf die Modelle schauen (\cref{sec:why,sec:explain}).
  \item Eine Datenqualitätsanalyse der 25 Hunderassen sowie dokumentierter, reproduzierbarer Code
        (\url{https://github.com/hsbi-dlcv-eskeif-dawalibi/hunderassen-vit}).
\end{itemize}

% Stand der Forschung: 3/4-1 Seite, mindestens 5 wissenschaftliche Quellen, nur relevante Arbeiten (Folie 16).
\section{Stand der Forschung}
\label{sec:related}

\paragraph{Vision Transformer und CNNs.}
Der ViT~\cite{dosovitskiy2021vit} zerlegt ein Bild in $16\times16$-Patches und verarbeitet sie mit einem
Transformer-Encoder~\cite{vaswani2017attention}. ViTs brauchen ein großes Pretraining und profitieren
bei wenigen Daten von Augmentierung und Regularisierung~\cite{steiner2022augreg}.
ResNet~\cite{he2016resnet} führte Residual-Blöcke ein; ConvNeXt~\cite{liu2022convnext} modernisiert
ResNet mit Designentscheidungen der Transformer und erreicht eine vergleichbare Accuracy.
Faltungen bringen Lokalität und Translationsäquivarianz mit, eine induktive Verzerrung, die bei wenigen
Daten hilft; globale Self-Attention muss die räumliche Struktur erst aus den Daten lernen, weshalb ViTs
stärker auf großes Pretraining angewiesen sind~\cite{dosovitskiy2021vit,steiner2022augreg}.

\paragraph{Transfer Learning und Linear Probing.}
Kornblith \etal~\cite{kornblith2019transfer} vergleichen logistische Regression auf eingefrorenen
Merkmalen mit vollem Fine-Tuning auf zwölf Datensätzen, darunter Oxford-IIIT Pet. Fine-Tuning hilft vor
allem bei feinkörnigen Datensätzen, die weit von ImageNet entfernt sind, etwa Autos und Flugzeuge; bei
Pets ist der Gewinn klein. Kumar \etal~\cite{kumar2022finetuning} zeigen, dass volles Fine-Tuning
vortrainierte Merkmale verzerren kann, was vor allem außerhalb der Trainingsverteilung schadet, während
Linear Probing sie erhält.

\paragraph{Parameter-effizientes Fine-Tuning.}
LoRA~\cite{hu2022lora} friert die vortrainierten Gewichte ein und lernt Änderungen niedrigen Rangs
$\Delta W = BA$ ausgewählter Gewichtsmatrizen; nach dem Training wird die Änderung in $W$ eingerechnet,
sodass die Inferenz nichts zusätzlich kostet. Hu \etal\ berichten, dass es genügt, nur die Query- und
Value-Projektionen anzupassen, dass schon ein sehr kleiner Rang ausreicht und dass mehrere Gewichtsarten
mit kleinerem Rang besser sind als eine einzige mit größerem Rang.

\paragraph{Feinkörnige Klassifikation von Hunderassen.}
Oxford-IIIT Pet~\cite{parkhi2012pets} (37 Katzen- und Hunderassen mit Segmentierungsmasken) und Stanford
Dogs~\cite{khosla2011dogs} (120 Rassen) sind Standard-Benchmarks für die Rassenerkennung, bei der sich
Klassen in feinen Details von Form und Fell unterscheiden.

\paragraph{Datenaugmentierung.}
RandAugment~\cite{cubuk2020randaugment} wendet zufällig gewählte Bildoperationen mit einer globalen
Stärke an; Random Erasing~\cite{zhong2020erasing} (ähnlich CutOut) maskiert zufällige Bildbereiche.
RandAugment ist Teil des AugReg-Rezepts für ViTs bei begrenzten Daten~\cite{steiner2022augreg}.

% Methode: 1-1 1/2 Seiten: Datensatz, Modelle, Parameter (Folie 16).
\section{Methode}
\label{sec:method}

\subsection{Datensatz und Datenqualität}
Wir verwenden die 25 Hunderassen von Oxford-IIIT Pet~\cite{parkhi2012pets} (CC BY-SA 4.0) und lassen
die 12 Katzenrassen weg. Eine einzige Tierart macht die Aufgabe feinkörnig; haben wir Hunde gewählt, weil es
etwa doppelt so viele Bilder wie von Katzen gibt. Mit 25 Klassen liegen wir über den vom Modul empfohlenen 5--10
Klassen, jede Rasse hat aber noch etwa 200 Bilder. Der offizielle \texttt{trainval}-Split wird in
Training und Validation geteilt (10\,\% pro Klasse, stratifiziert, fester Split-Seed); der offizielle
\texttt{test}-Split wird einmal nach der Modellauswahl verwendet (\cref{tab:splits}).

\begin{table}[t]
  \centering
  \caption{Aufteilung der 25 Hunderassen in Training, Validation und Test.}
  \label{tab:splits}
  \small
  \begin{tabular}{lrl}
    \toprule
    Split & Bilder & Quelle \\
    \midrule
    Training   & 2\,242 & \texttt{trainval} ohne Validation \\
    Validation & 250    & 10\,\% pro Rasse, stratifiziert \\
    Test       & 2\,486 & offizieller \texttt{test}-Split \\
    \bottomrule
  \end{tabular}
\end{table}

\paragraph{Datenqualität.}
Wir haben alle 4\,978 Bilder geprüft. Die Klassen sind nahezu balanciert
(189--200 Bilder pro Rasse). Die Bildgröße liegt im Median bei $500\times375$\,px; 133 Bilder haben eine kürzere Seite unter
224\,px und werden hochskaliert. Der Hund bedeckt im Mittel 40\,\% des Bildes, in 123 Bildern weniger als
10\,\%. Exakte Duplikate gibt es nicht, daher wurde kein Bild entfernt. \Cref{fig:examples} zeigt ein
Beispiel pro Rasse.

\paragraph{Augmentierung.}
Unser Standard (\emph{light}) nutzt zufällige Ausschnitte (60--100\,\% des Bildes) und horizontales
Spiegeln; vertikales Spiegeln und Farbänderungen vermeiden wir, weil Orientierung und Fellfarbe
Rassemerkmale sind. \emph{strong} ergänzt RandAugment~\cite{cubuk2020randaugment} (zufällig gewählte Bildoperationen) und Random
Erasing~\cite{zhong2020erasing} (zufällig maskierte Bildbereiche), \emph{none} nutzt nur Skalierung und zentralen Zuschnitt.

\begin{figure}[t]
  \centering
  \resultfigure{paper/beispiele.png}{\linewidth}{Beispiel pro Rasse mit Maskenumriss}
  \caption{Je ein Trainingsbild der 25 Hunderassen (das erste der Rasse im Trainingssplit). Gelb: Umriss der
  Segmentierungsmaske des Tieres. Manche Rassen sehen sich sehr ähnlich, etwa American Pit Bull Terrier und
  Staffordshire Bull Terrier.}
  \label{fig:examples}
\end{figure}

\subsection{Modelle und Training}
Alle Modelle sind auf ImageNet~\cite{deng2009imagenet,russakovsky2015imagenet} vortrainiert und werden mit Hugging Face
Transformers~\cite{wolf2020transformers} auf PyTorch~\cite{paszke2019pytorch} geladen: ViT-B/16
(\texttt{google/vit-base-patch16-224}, auf ImageNet-21k vortrainiert und auf ImageNet-1k feinjustiert; 86\,M Parameter mit unserem Kopf für 25 Klassen), ResNet-50
(\texttt{microsoft/resnet-50}, 24\,M) und ConvNeXt-T (\texttt{facebook/convnext-tiny-224}, 28\,M). Wir
trainieren mit AdamW~\cite{loshchilov2019adamw} (Weight Decay 0,05) und Cosine-Schedule~\cite{loshchilov2017sgdr}
mit 10\,\% Warm-up, bis zu 10 Epochen mit
Early Stopping, Batchgröße 32 und $224\times224$ Pixeln Eingabe. Jede Konfiguration läuft mit drei Seeds;
wir berichten Mittelwert $\pm$ Standardabweichung auf dem Testsplit. Der Checkpoint wird auf dem
Validierungssplit ausgewählt.

\paragraph{Auswertung.}
Wir berichten Accuracy und Macro-F1. Einen Unterschied zwischen zwei Methoden interpretieren wir nur, wenn
die Differenz der Mittelwerte größer ist als das Doppelte der größeren Seed-Standardabweichung beider
Methoden (\emph{Schwelle}); alle Aussagen wie \glqq innerhalb der Streuung\grqq{} beziehen sich darauf. Mit drei
Seeds ist das ein grobes Kriterium und kein statistischer Test; die Schwelle ist konservativer als ein
gepaarter Vergleich je Seed (siehe Limitationen).

\subsection{Anpassungsmethoden}
\label{sec:methods}
\paragraph{Linear Probing.} Der ViT-Backbone ist eingefroren; trainiert wird nur der neue lineare
Klassifikator auf dem \texttt{[CLS]}-Token (Lernrate $10^{-3}$).

\paragraph{LoRA.} Für ein eingefrorenes vortrainiertes Gewicht $W\in\mathbb{R}^{d_\text{out}\times
d_\text{in}}$ lernt LoRA~\cite{hu2022lora} $A\in\mathbb{R}^{r\times d_\text{in}}$ und
$B\in\mathbb{R}^{d_\text{out}\times r}$ mit Rang $r\ll d$:
\begin{equation}
  h = W x + \frac{\alpha}{r}\, B A x .
  \label{eq:lora}
\end{equation}
In \cref{eq:lora} wird $B$ mit Nullen initialisiert, sodass das Training beim vortrainierten Modell
beginnt. Wir passen die Query- und Value-Projektionen aller 12 Blöcke mit $r\in\{4,8,16\}$ an und für
$r=8$ zusätzlich alle linearen Layer jedes Blocks (Key, Attention-Ausgabe, beide MLP-Layer). Wir setzen
$\alpha=2r$, verwenden Dropout 0,1 im LoRA-Pfad und trainieren zusätzlich den Klassifikationskopf (Lernrate
$5\cdot10^{-4}$). Nach dem Training wird $\frac{\alpha}{r}BA$ zu $W$ addiert; das angepasste Modell hat
damit dieselbe Architektur und dieselben FLOPs wie der ursprüngliche ViT.

\paragraph{Volles Fine-Tuning.} Alle Parameter werden trainiert (Lernrate $5\cdot10^{-5}$); das ist die
Referenz, mit der die effizienten Methoden verglichen werden. Auch die CNNs werden vollständig
angepasst. Gemäß Vorgabe des Dozenten werden sie nicht von Grund auf trainiert, sondern starten mit vortrainierten
ImageNet-Gewichten und werden nur an unsere Aufgabe angepasst (Transfer Learning); das Training ohne
Pretraining in \cref{sec:why} dient ausschließlich als Vergleich.

\paragraph{Lernraten.} Die Lernraten der ViT-Methoden wurden einmal pro Methode festgelegt und nicht
optimiert. Unsere ersten CNN-Läufe verwendeten die
ViT-Lernrate $5\cdot10^{-5}$ und waren deutlich untertrainiert (Validation-Loss 1,35 bei ResNet-50). Daher
haben wir die CNN-Lernrate auf dem \emph{Validierungssplit} aus $\{5\cdot10^{-5}, 2\cdot10^{-4},
5\cdot10^{-4}, 10^{-3}\}$ mit je drei Seeds gewählt; $2\cdot10^{-4}$ war für beide CNNs am besten
(\cref{tab:lr}).

\begin{table}[h]
  \centering
  \caption{Wahl der CNN-Lernrate auf dem Validierungssplit (Mittelwert $\pm$ Std über 3 Seeds; Loss:
  Cross-Entropy; fett: gewählte Lernrate).}
  \label{tab:lr}
  \small
  \resizebox{\linewidth}{!}{\begin{tabular}{lcccc}
    \toprule
    & \multicolumn{2}{c}{ResNet-50} & \multicolumn{2}{c}{ConvNeXt-T} \\
    Lernrate & Val.-Acc.\ (\%) & Val.-Loss & Val.-Acc.\ (\%) & Val.-Loss \\
    \midrule
    $5\cdot10^{-5}$ & $91{,}47\pm1{,}29$ & 1,349 & $96{,}40\pm0{,}69$ & 0,303 \\
    $2\cdot10^{-4}$ & $\mathbf{97{,}47\pm0{,}46}$ & 0,146 & $\mathbf{97{,}47\pm0{,}46}$ & 0,112 \\
    $5\cdot10^{-4}$ & $96{,}40\pm1{,}06$ & 0,110 & $96{,}13\pm0{,}83$ & 0,157 \\
    $10^{-3}$ & $96{,}13\pm1{,}80$ & 0,154 & $94{,}80\pm0{,}80$ & 0,190 \\
    \bottomrule
  \end{tabular}}
\end{table}

% Ergebnisse: 1-1 1/2 Seiten (Folie 16). Nur Werte aus results/ms1/REPORT.md, Mittelwert +- Std über 3 Seeds.
\section{Ergebnisse}
\label{sec:results}

\paragraph{Vergleich aller Methoden.}
\Cref{tab:methods} und \cref{fig:efficiency} zeigen alle Methoden. LoRA mit $r=8$ erreicht die
höchste Test-Accuracy ($96{,}21\pm0{,}24$\,\%), gefolgt von LoRA $r=4$ ($96{,}11\pm0{,}21$\,\%) und
Linear Probing ($96{,}03\pm0{,}12$\,\%), das nur 0,02\,\% der Parameter trainiert. Die Unterschiede
zwischen diesen drei liegen innerhalb der Seed-Streuung. Volles Fine-Tuning ($95{,}39\pm0{,}39$\,\%)
liegt 0,82 Prozentpunkte (pp) unter der besten LoRA-Variante, mehr als unsere Schwelle. Die CNNs
erreichen $94{,}78\pm0{,}46$\,\% (ConvNeXt-T) und $94{,}22\pm0{,}85$\,\% (ResNet-50); die beste
ViT-Methode liegt 1,42\,pp über dem besten CNN. Bei der Inferenz braucht der ViT-B/16 jedoch 35,1\,GFLOPs
pro Bild, ResNet-50 8,2 und ConvNeXt-T 8,9.

\begin{table*}[t]
  \centering
  \caption{Testergebnisse (Mittelwert $\pm$ Standardabweichung über 3 Seeds; fett: bester Wert).
  Trainierbar: im Training aktualisierte Parameter. Zeit: Training eines Laufs auf einer RTX~5090. GFLOPs: Rechenaufwand
  pro Bild bei der Inferenz. LoRA-Änderungen werden nach dem Training eingerechnet (keine zusätzlichen
  Inferenzkosten).}
  \label{tab:methods}
  \small
  \begin{tabular}{lccccc}
    \toprule
    Methode & Acc.\ (\%) & Macro-F1 (\%) & Trainierbar & Zeit (min) & GFLOPs \\
    \midrule
    \methodrows
    \bottomrule
  \end{tabular}
\end{table*}

\begin{figure}[t]
  \centering
  \resultfigure{paper/effizienz.png}{\linewidth}{Accuracy über trainierbaren Parametern}
  \caption{Test-Accuracy über der Zahl trainierbarer Parameter (logarithmische Skala). Punkte: Mittelwert über
  3 Seeds, Balken: $\pm$ eine Standardabweichung. Mit weniger trainierbaren Parametern bleibt die Accuracy
  gleich oder höher; volles Fine-Tuning und die CNNs liegen darunter.}
  \label{fig:efficiency}
\end{figure}

\paragraph{LoRA-Parameter.}
Mit Anpassung von Query und Value beträgt die Accuracy 96,11, 96,21 und 95,96\,\% für $r=4$, 8 und 16
(0,17, 0,31 und 0,61\,M trainierbare Parameter). Die Anpassung aller linearen Layer mit $r=8$
(1,35\,M Parameter) ergibt $95{,}82\pm0{,}18$\,\%, 0,39\,pp unter Query/Value allein und damit innerhalb
der Seed-Streuung.

\paragraph{Trainingsverlauf.}
Die beste Validation-Accuracy erreicht Linear Probing nach $3{,}7\pm1{,}2$ Epochen, die übrigen Methoden
nach 5--7 (\cref{fig:curves}).

\begin{figure}[h]
  \centering
  \resultfigure{paper/lernkurven.png}{\linewidth}{Lernkurven}
  \caption{Validation-Accuracy je Epoche (Mittelwert über 3 Seeds; die y-Achse beginnt bei 68\,\%). Linear
  Probing liegt schon nach einer Epoche bei 94\,\%. Die CNNs starten niedriger (ResNet-50: 70\,\%) und
  bleiben nach dem Aufholen unter den LoRA-Varianten.}
  \label{fig:curves}
\end{figure}

\paragraph{Augmentierung.}
Für LoRA $r=8$ ergibt \emph{light} $96{,}21\pm0{,}24$\,\%, \emph{none} $95{,}99\pm0{,}16$\,\% und
\emph{strong} $95{,}91\pm0{,}36$\,\%; die Unterschiede liegen innerhalb der Seed-Streuung.

\paragraph{Fehler.}
Für LoRA $r=4$ (auf Validation gewählt) sind American Pit Bull Terrier (F1 77,2\,\%), Staffordshire
Bull Terrier (78,2\,\%) und American Bulldog (88,7\,\%) am schwierigsten (Anhang~\ref{sec:perbreed}).
Häufigste Verwechslung: Pit Bull $\rightarrow$ Staffordshire (62 Bilder).

% Diskussion: 1-1 1/2 Seiten (Folie 16). Dem Dozenten besonders wichtig: Was bedeuten die Ergebnisse?
% Warum ist eine Methode besser? Auf welchen Daten funktioniert eine Variante?
% Alle Zahlen aus results/ms1/REPORT.md und results/why/REPORT.md. Regel: Ein Abstand zählt nur, wenn er
% größer ist als die doppelte größere Seed-Standardabweichung der beiden Methoden.
\section{Diskussion}
\label{sec:discussion}

\paragraph{Welche Methode ist die beste, und warum?}
LoRA $r=8$ ist nominell die beste Methode (96,21\,\%), doch Linear Probing (96,03\,\%) und LoRA $r=4$
(96,11\,\%) sind innerhalb der Seed-Streuung nicht von ihr zu unterscheiden. Nur volles Fine-Tuning
($-0{,}82$\,pp) und die CNNs ($-1{,}42$ und $-1{,}99$\,pp) sind klar schlechter. Die Zusatzexperimente
(\cref{sec:why}) legen nahe, warum ein eingefrorener Backbone genügt: Die vortrainierten Merkmale
trennen die Rassen bereits. Eine logistische Regression auf eingefrorenen ViT-B/16-Merkmalen erreicht
95,9\,\%, mit nur einem Bild pro Rasse schon 84,7\,\%. Zudem sind 22 unserer 25 Rassen
ImageNet-1k-Klassen~\cite{russakovsky2015imagenet}; bei Auswahl unter diesen 22 Rassen erreichen die
ursprünglichen ImageNet-Köpfe 96,2--96,8\,\%, ohne ein einziges Bild unseres Trainingssatzes gesehen zu
haben. Die Anpassung muss also keine Rassenmerkmale lernen, sondern vor allem vorhandene Merkmale auf
unsere Klassen abbilden, und das leistet ein linearer Kopf. Volles Fine-Tuning aktualisiert dagegen 86\,M
Parameter mit 2\,242 Bildern und landet leicht unter den effizienten Methoden. Das passt zur Beobachtung,
dass volles Fine-Tuning gute vortrainierte Merkmale verzerren kann~\cite{kumar2022finetuning}, auch wenn
Kumar \etal\ den Schaden vor allem außerhalb der Trainingsverteilung feststellen. Bei ResNet-50 zeigt sich
ein ähnliches Muster: Eingefrorene Merkmale mit Linear Probe erreichen 96,1\,\%, das vollständig
feinjustierte Netz 94,2\,\% (ein Lauf der Probe gegenüber drei Fine-Tuning-Seeds und ein anderes
Trainingsverfahren, der Vergleich ist daher nur ein Hinweis).

\paragraph{Lohnt sich LoRA? (A1, A2)}
A1 ist bestätigt: Mit 0,19--0,70\,\% trainierbaren Parametern ist LoRA nie schlechter als volles
Fine-Tuning, und LoRA $r=8$ ist sogar signifikant besser ($+0{,}82$\,pp bei einer Schwelle von
0,78\,pp); die anderen Ränge liegen innerhalb der Streuung. A2 ist widerlegt: Linear Probing liegt nur
0,17\,pp hinter der besten LoRA-Variante, innerhalb des Rauschens, trainiert dabei 16-mal weniger
Parameter als LoRA $r=8$ und ist am schnellsten fertig (im Mittel 0,3\,min; seine beste Epoche,
$3{,}7\pm1{,}2$, ist nominell auch die früheste). Auf diesem Datensatz gilt die Empfehlung der Vorlesung,
bei weniger als 10\,k Bildern nur den Kopf zu trainieren, sogar für feinkörnige Rassen. Gegenüber
Adapter-Methoden, die zur Inferenz zusätzliche Module brauchen, hat LoRA den Vorteil, dass die Änderung in
die Gewichte eingerechnet wird.

\paragraph{Rang und angepasste Layer (A3).}
Der erste Teil von A3 ist bestätigt, der zweite nicht. Die Ränge 4, 8 und 16 unterscheiden sich um etwa
0,25\,pp, und die Anpassung aller linearen Layer mit $r=8$ (1,35\,M Parameter) ist nicht besser als nur
Query/Value ($-0{,}39$\,pp, innerhalb des Rauschens). In Analogie zum niedrigen intrinsischen Rang der
nötigen Änderung, den Hu \etal~\cite{hu2022lora} für Sprachmodelle berichten, scheint zusätzliche
Kapazität im Adapter ungenutzt zu bleiben, wenn die eingefrorenen Merkmale bereits gut sind.

\paragraph{ViT gegen CNN (A4).}
Das beste CNN (ConvNeXt-T, 94,78\,\%) liegt 1,42\,pp unter dem besten ViT. Der Abstand ist kleiner als
2\,pp, übersteigt aber die Seed-Streuung (Schwelle 0,92\,pp); A4 ist daher im strengen Sinn widerlegt:
Die Accuracies liegen nah beieinander, unterscheiden sich aber systematisch. Unsere Probes auf
eingefrorenen Merkmalen deuten auf die Ursache. Mit Linear Probe sind die Merkmale von ResNet-50
(96,1\,\%, nur ImageNet-1k) so gut wie die von ViT-B/16 (95,9\,\%, ImageNet-21k und 1k). Der Abstand
scheint also von der Anpassung der CNNs durch volles Fine-Tuning zu kommen, das auch beim ViT die
schwächste Methode ist, und nicht von einem besseren Pretraining des ViT. Eine methodische Lehre betrifft
die Lernrate: Mit der ViT-Lernrate ($5\cdot10^{-5}$) war ResNet-50 untertrainiert (Validation-Accuracy
91,5\,\%, Validation-Loss 1,35); erst eine pro Architektur auf dem Validierungssplit gewählte Lernrate
macht den Vergleich fair (\cref{tab:lr}).

\paragraph{Auf welchen Daten würde eine andere Methode gewinnen?}
Unser Ergebnis hängt davon ab, wie nah die Zielaufgabe an den Pretraining-Daten liegt. Die drei Rassen
ohne ImageNet-Klasse erreichen mit Linear Probing einen etwas kleineren F1-Wert (93,8 gegenüber
96,2\,\%), vor allem wegen des American Bulldog, der mit ähnlichen Rassen verwechselt wird. Für Daten fern
von ImageNet (andere Bilddomänen, Klassen ohne Entsprechung im Pretraining) oder für deutlich größere
Datensätze dürften eingefrorene Merkmale nicht mehr genügen, und LoRA oder volles Fine-Tuning sollten
vorn liegen, wie Kornblith \etal~\cite{kornblith2019transfer} für Autos und Flugzeuge beobachten und wie
es die Strategietabelle der Vorlesung nahelegt. Training ohne Pretraining ist bei dieser Datenmenge keine
Alternative: ViT-B/16 erreicht nur 11,2\,\% und ResNet-50 31,1\,\%; der Vorsprung des CNN geht
vermutlich auf die induktive Verzerrung der Faltungen zurück.

\paragraph{Augmentierung und Datenqualität.}
\emph{none}, \emph{light} und \emph{strong} unterscheiden sich um höchstens 0,3\,pp, innerhalb des
Rauschens. Bei so guten Merkmalen und einem eingefrorenen oder niedrigrangig angepassten Backbone gibt es
wenig Overfitting, das Augmentierung verhindern könnte, und starke Augmentierung kann Rassemerkmale wie
die Fellfarbe verändern. Die Daten sind balanciert und enthalten keine exakten Duplikate. 8,6\,\% der
Testbilder haben ein sehr ähnliches Trainingsbild (CLIP-Ähnlichkeit $\geq0{,}95$); ohne sie sinkt die
Accuracy um 0,34\,pp, die Ergebnisse sind durch Fast-Duplikate also höchstens leicht geschönt. Die
verbleibenden Fehler konzentrieren sich: Die vier Bully-Rassen machen 16\,\% der Testbilder aus,
verursachen aber 59\,\% aller Fehler. Diese Rassen sehen sich auf Fotos sehr ähnlich und unterscheiden
sich vor allem in Größe und Proportionen, die ein Bild ohne Maßstab kaum zeigt.

\paragraph{Limitationen.}
Wir verwenden einen Datensatz, dessen Klassen großteils in ImageNet vorkommen; der Vorteil von Linear
Probing muss sich daher nicht auf andere Domänen übertragen. Nur die CNN-Lernrate wurde optimiert, die
ViT-Methoden nutzen je eine feste Lernrate. Drei Seeds und ein Validierungssplit mit 250 Bildern (ein
Bild entspricht 0,4\,pp; die Modellauswahl für Station~2 war ein Gleichstand) begrenzen die Auflösung
kleiner Unterschiede. Die Trainingszeiten wurden auf einer einzigen GPU gemessen.

\paragraph{Fazit und Ausblick.}
Für einen kleinen feinkörnigen Datensatz nahe an ImageNet erreicht ein neuer linearer Kopf auf einem
eingefrorenen ViT-B/16 96\,\% mit 0,02\,\% trainierbaren Parametern und dem kürzesten Training; LoRA
bringt keinen messbaren Gewinn, und volles Fine-Tuning ist leicht schlechter ($-0{,}8$\,pp). Die auf dem
Validierungssplit gewählte Methode, LoRA $r=4$, dient in Station~2 als Teacher, den wir in deutlich
kleinere ViTs destillieren.

{
    \small
    \bibliographystyle{ieeenat_fullname}
    \bibliography{main}
}

% Appendix / supplementary material (does not count towards the 6 pages), in the
% format of the author kit (sec/X_suppl.tex). Remove the line if not needed.
% Anhang / Zusatzmaterial: zählt nicht zu den 6 Seiten (Folie 16). Format wie im Author-Kit
% (sec/X_suppl.tex), Überschrift auf Deutsch. Zahlen aus results/ms1, results/why, results/explain.
\clearpage
\setcounter{page}{1}
\twocolumn[\centering\Large\textbf{\thetitle}\\\vspace{0.5em}Anhang\\\vspace{1.0em}]
\appendix

\section{Hyperparameter}
\label{sec:hyperparameters}
\Cref{tab:hyper} listet alle Konfigurationen (\texttt{configs/ms1/*.yaml}). Gemeinsam für alle: AdamW,
Weight Decay 0,05, Cosine-Schedule mit 10\,\% Warm-up, bis zu 10 Epochen, Early Stopping (Geduld 3),
Batchgröße 32, $224\times224$ Pixel Eingabe, Augmentierung \emph{light}, Seeds 1--3.

\begin{table}[h]
  \centering
  \caption{Lernrate und trainierbare Parameter je Methode.}
  \label{tab:hyper}
  \small
  \resizebox{\linewidth}{!}{\begin{tabular}{lll}
    \toprule
    Methode & Lernrate & Trainierbar \\
    \midrule
    ViT Linear Probing & $10^{-3}$ & Kopf \\
    ViT LoRA $r\in\{4,8,16\}$, $\alpha=2r$, Dropout 0,1, Query/Value & $5\cdot10^{-4}$ & LoRA + Kopf \\
    ViT LoRA $r=8$, alle linearen Layer & $5\cdot10^{-4}$ & LoRA + Kopf \\
    ViT volles Fine-Tuning & $5\cdot10^{-5}$ & alle \\
    ResNet-50, ConvNeXt-T & $2\cdot10^{-4}$ (Validation) & alle \\
    \bottomrule
  \end{tabular}}
\end{table}

\section{Warum sind die Accuracies so hoch?}
\label{sec:why}
Alle Zahlen aus \texttt{results/why/REPORT.md} (\texttt{make why}).

\paragraph{Pretraining.} Zufällig initialisiert und mit denselben Daten und derselben Auswertung
trainiert (30 Epochen, größere Lernrate), erreicht ViT-B/16 $11{,}18\pm2{,}61$\,\% und ResNet-50
$31{,}07\pm0{,}41$\,\% Test-Accuracy, gegenüber $95{,}39\pm0{,}39$ und $94{,}22\pm0{,}85$\,\% mit
ImageNet-Gewichten (Zufall: 4\,\%).

\paragraph{Vorwissen.} 22 der 25 Rassen sind ImageNet-1k-Klassen~\cite{russakovsky2015imagenet}; es
fehlen nur American Bulldog, Havanese und Shiba Inu. Ganz ohne Training klassifizieren die
ursprünglichen ImageNet-Köpfe unserer Checkpoints die Testbilder dieser 22 Rassen (Auswahl unter den 22)
mit 96,20\,\% (ViT-B/16), 96,75\,\% (ResNet-50) und 96,52\,\% (ConvNeXt-T). LoRA $r=4$ erreicht auf
denselben Bildern 95,88\,\%, bei Auswahl unter allen 25 Rassen. Die Zuordnung von Rasse zu
ImageNet-Klasse erfolgt über den Namen (Tabelle in \texttt{animal\_vit/why.py}); sie ist eine Annahme. Der
American Pit Bull Terrier wird der Klasse ``American Staffordshire terrier'' zugeordnet, Japanese Chin
``Japanese spaniel'', Scottish Terrier ``Scotch terrier'' und German Shorthaired ``German short-haired
pointer''. Bei diesen Rassen kann die Zuordnung strittig sein.

\paragraph{Eingefrorene Merkmale.} Eine logistische Regression~\cite{pedregosa2011sklearn} auf eingefrorenen Merkmalen (alle
Trainingsbilder) erreicht 95,94\,\% (ViT-B/16) und 96,06\,\% (ResNet-50), gegenüber 95,39 und 94,22\,\%
nach vollem Fine-Tuning. Mit einem Bild pro Rasse erreicht sie $84{,}66\pm2{,}63$ und
$87{,}41\pm1{,}91$\,\% (\cref{fig:fewshot} im Hauptteil).

\paragraph{Fast-Duplikate.} Für jedes Testbild haben wir die Kosinus-Ähnlichkeit zum ähnlichsten Trainings-
oder Validierungsbild mit Embeddings von CLIP ViT-B/16~\cite{radford2021clip} berechnet. 215 Testbilder
(8,6\,\%) haben einen Nachbarn mit Ähnlichkeit $\geq0{,}95$; ohne sie ändert sich die Accuracy von LoRA
$r=4$ von 96,11 auf 95,77\,\%.

\paragraph{Rest-Fehler.} Die vier Bully-Rassen (American Bulldog, American Pit Bull Terrier,
Staffordshire Bull Terrier, Boxer) machen 16\,\% der Testbilder und 59\,\% aller Fehler von LoRA $r=4$
aus (Accuracy 84,9\,\% gegenüber 98,2\,\% bei den übrigen Rassen). 71 Testbilder sind in allen drei
Seeds falsch klassifiziert.


\section{Warum diese Rasse und nicht die andere?}
\label{sec:explain}
Erklärungen für LoRA $r=4$, ConvNeXt-T und ResNet-50 (Seed 1; \texttt{results/explain/REPORT.md},
\texttt{make explain}). Wir nutzen Grad-CAM~\cite{selvaraju2017gradcam} (Implementierung:
pytorch-grad-cam~\cite{gildenblat2021gradcam}), Occlusion~\cite{zeiler2014visualizing} (Implementierung:
Captum~\cite{kokhlikyan2020captum}) und Attention Rollout~\cite{abnar2020rollout} (eigener Code). Ein
\emph{kontrastives} Grad-CAM verwendet das Ziel $\text{logit}_A-\text{logit}_B$ und hebt Evidenz für
Rasse $A$ und gegen Rasse $B$ hervor (vgl.~\cite{prabhushankar2020contrastive}). In einem
Deletion-Test~\cite{petsiuk2018rise} werden die auffälligsten $16\times16$-Patches zuerst grau
übermalt; Occlusion senkt die Klassenwahrscheinlichkeit beim ViT und bei ConvNeXt-T am schnellsten
(Fläche unter der Kurve 0,64 und 0,59 gegenüber 0,92 und 0,89 beim zufälligen Löschen). Gemessen mit der
jeweils treuesten Methode liegt die Saliency in der annotierten Kopf-Box (Validierungssplit) bei allen
drei Modellen beim 3,1- bis 3,7-Fachen ihres Flächenanteils. Grad-CAM liefert für den ViT einen deutlich
kleineren Wert (1,3), passend zu bekannten Grenzen von Grad-CAM bei Transformern~\cite{chefer2021transformer}.
\Cref{fig:head} im Hauptteil zeigt Beispielkarten mit den Kopf-Boxen, \cref{fig:contrastive} kontrastive Karten für
die häufigsten Verwechslungen.




\begin{figure}[t]
  \centering
  \resultfigure{explain/contrastive_convnext_t.png}{\linewidth}{Kontrastive Karten}
  \caption{Warum Rasse $A$ und nicht $B$? Für die häufigsten Verwechslungen: Grad-CAM für $A$ und $B$,
  kontrastive Karte $A-B$ und Occlusion für $A$ (ConvNeXt-T, Testbilder).}
  \label{fig:contrastive}
\end{figure}

\section{Software und KI-Hilfsmittel}
\label{sec:tools}
Die Experimente nutzen PyTorch~\cite{paszke2019pytorch}, Hugging Face
Transformers~\cite{wolf2020transformers}, scikit-learn~\cite{pedregosa2011sklearn}, pytorch-grad-cam~\cite{gildenblat2021gradcam}
und Captum~\cite{kokhlikyan2020captum}.
% Erklärung zur KI-Nutzung (Modulregeln, Vorlesung 01, Folie 13); in der Präsentation mit dem Dozenten besprochen.
Bei Code, Auswertung und Formulierung haben wir einen KI-Assistenten eingesetzt, wie in der Präsentation mit
dem Dozenten besprochen. Forschungsfragen, Experimente und die Prüfung aller Ergebnisse stammen von uns;
alle Zahlen stammen aus den protokollierten Experimenten (\texttt{results/}).


\end{document}
